% authority: science_draft
\documentclass[11pt]{article}

\usepackage[margin=1in]{geometry}
\usepackage{amsmath,amssymb}
\usepackage{booktabs}
\usepackage{longtable}
\usepackage{enumitem}

\newcommand{\AuditVerdict}{\mathsf{source\_pure\_as\_written\_pending\_refuter\_stress}}
\newcommand{\EqMS}{\mathrm{EqMS}_{\mathrm{src}}^{\mathrm{cert}}}
\newcommand{\CTr}{C_{\mathrm{transport}}^{\mathrm{src}}}
\newcommand{\CInv}{C_{\mathrm{invariance}}^{\mathrm{src}}}
\newcommand{\CFact}{C_{\mathrm{factorization}}^{\mathrm{src}}}
\newcommand{\CNTI}{C_{\mathrm{no\text{-}target}}^{\mathrm{src}}}
\newcommand{\RREObj}{\mathrm{RR}_{E}^{\mathrm{decl}}}
\newcommand{\MetricData}{\mathrm{MetricData}}
\newcommand{\geff}{g_{\mathrm{eff}}}

\title{Matter-Semantics Equivalence Theorem Smuggling Audit v1}
\author{Smuggling Auditor Packet RT-20260702-059}
\date{2026-07-02}

\begin{document}
\maketitle

\section{Control Status}

This document implements v15 P2-T04. It audits the P2-T03 artifact
\emph{Narrow source-side matter-semantics equivalence theorem attempt v1}
from \texttt{RT-20260702-058}.

This is a draft/control source-purity audit. It is not a canonical ontology
edit, not source-law adoption, not \(\mathrm{RR}_{E}\) transport-law
adoption, not matter-semantics adoption, not detector-semantics adoption, not
coupling-law adoption, not matter-coupling derivation, not stress-energy
semantics, not a matter action, not Einstein equations, not benchmark
promotion, and not completed derivation.

\section{Audited Object}

The audited theorem states that declared source-side matter-semantics objects
\(A_{\mathrm{src}}\) and \(B_{\mathrm{src}}\) satisfy
\[
  A_{\mathrm{src}} \EqMS B_{\mathrm{src}}
\]
when at least one explicit source certificate
\[
  \CTr(A,B),\qquad \CInv(A,B),\qquad \CFact(A,B;F)
\]
is supplied, the no-target certificate \(\CNTI\) passes, declared-object
indexing is consistent, and fail-closed branches remain inactive. The theorem
also records certificate-gap and malformed/imported-certificate fail-closed
branches.

\section{Audit Criteria}

The audit checks whether the theorem obtains its result by importing any
forbidden authority:

\begin{enumerate}[label=(\arabic*)]
  \item target manifold, target topology, target atlas, target differentiable
        structure, target Lorentzian signature, or target metric;
  \item empirical clocks, rods, detector protocols, detector semantics, or
        empirical observer readout semantics;
  \item benchmark behavior or exact-GR success;
  \item \(\MetricData(E)\), \(\geff\), stress-energy tensor, matter action,
        or Einstein equations;
  \item validator PASS, registry row, Gate Chair status outside exact scope,
        generated derivative, file order, commit state, role identity,
        handoff, approval, local cache, or other process authority as proof;
  \item scoped evidence/precondition status as adoption or as a substitute for
        explicit source certificates;
  \item conditional theorem status as matter semantics, detector semantics,
        coupling law, matter coupling, benchmark recovery, or completed
        derivation.
\end{enumerate}

\section{Findings}

\subsection{No target-manifold, topology, atlas, signature, or metric import}

The theorem's positive proof proceeds by cases on explicit source
certificates. The transport case uses source labels, guards, and record
indices. The invariance case uses source relabeling or admissible source
variation. The factorization case uses a declared source object \(F\) and
object-indexed \(\RREObj(F)\). The proof does not require target topology,
target smooth atlas, target differentiability, target Lorentzian signature,
target metric, proper time, empirical clocks, or rods.

\subsection{No detector-semantics import detected as written}

The theorem keeps source readout or response tokens inside source-side syntax.
It explicitly denies detector equivalence, physical matter equivalence,
empirical equivalence, and benchmark equivalence. No detector protocol or
empirical observer semantics is used as a mathematical premise.

\subsection{No stress-energy, matter-action, or Einstein-equation import}

The theorem neither imports nor constructs stress-energy semantics, a
stress-energy tensor, matter-action semantics, a matter action, or Einstein
equations. The Distance-to-GR section states that those burdens remain
unchanged and blocked.

\subsection{No process-authority laundering detected as written}

The theorem states that generated derivatives, registry metadata, validators,
roles, handoffs, approvals, local caches, file order, and commit state may not
be used as mathematical premises. Its proof uses the P2-T02 source-side
definition of \(\EqMS\), not validation or registry status, as the local
logical basis.

\subsection{No status laundering detected as written}

The theorem's positive status is conditional and scoped:
\[
  \texttt{conditional\_or\_scoped\_source\_side\_equivalence\_theorem\_candidate}.
\]
It is not rendered as bare accepted status. It does not adopt
\(\CTr\), \(\CInv\), \(\CFact\), \(\CNTI\), \(\mathrm{PositiveMSProfile}_{v1}\),
\(\mathrm{SourceMatterSemanticsAdoptionReadinessLaw}_{v1}\), or
\(\mathrm{RR}_{E}\mathrm{TransportCompletenessOrInvarianceLaw}_{v1}\) as
source laws.

\subsection{Scoped evidence is not used as certificate evidence}

The theorem explicitly says scoped evidence/precondition artifacts may be
cited only within their tracked scopes and do not instantiate a concrete pair
\((A_{\mathrm{src}},B_{\mathrm{src}})\) with a complete certificate bundle.
Therefore the theorem does not convert scoped evidence into adoption or
certificate witnesses.

\section{Verdict}

The audit verdict is
\[
  \AuditVerdict .
\]

No target import, detector-semantics import, metric import, stress-energy
import, process-authority laundering, or status laundering is detected in the
theorem packet as written.

This verdict makes the theorem packet ready for P2-T05 Refuter stress. It does
not prove any concrete certificate instance, does not remove the certificate
gap branch, does not adopt source laws or matter semantics, and does not
advance the Distance-to-GR ledger.

\section{Required P2-T05 Stress Targets}

The next Refuter packet should stress at least:

\begin{enumerate}[label=(\alph*)]
  \item deletion of \(\CTr\), \(\CInv\), and \(\CFact\) certificates;
  \item malformed source transport;
  \item source-invariance mismatch;
  \item source-factorization object mismatch;
  \item \(\RREObj(F)\) separation violation pressure;
  \item detector-semantics collapse pressure;
  \item target metric and proper-time import pressure;
  \item \(\CNTI\) misuse as positive semantics;
  \item finite/local-to-global overread;
  \item scoped evidence as adoption overread;
  \item source-extension as derivation overread;
  \item validator, registry, role, handoff, approval, generated-derivative,
        file-order, local-cache, commit, or checkpoint proof-laundering
        pressure.
\end{enumerate}

\section{Distance-to-GR Status}

The Distance-to-GR status is unchanged. The audit classifies the theorem
packet as source-pure as written pending stress. Matter coupling remains
undischarged. \(\MetricData(E)\), \(\geff\), stress-energy semantics, matter
action, Einstein equations, benchmark promotion, and completed derivation
remain blocked.

\section{Non-Conclusions}

This audit does not authorize canonical ontology edit, source-law adoption,
\(\mathrm{RR}_{E}\) source-law adoption,
\(\mathrm{RR}_{E}\mathrm{TransportCompletenessOrInvarianceLaw}_{v1}\)
adoption, unrestricted \(\mathrm{RR}_{E}\) theorem authority,
\(\mathrm{PositiveMSProfile}_{v1}\) adoption,
\(\mathrm{SourceMatterSemanticsAdoptionReadinessLaw}_{v1}\) adoption as law,
source-extension data adoption beyond exact scoped Gate Chair results,
\(\MetricData(E)\) adoption, \(\geff\) adoption or scope expansion,
coupling-law adoption, matter-semantics adoption, detector-semantics adoption,
matter-coupling derivation or adoption, stress-energy semantics,
stress-energy tensor, matter action, Einstein equations, benchmark promotion,
benchmark closure, completed derivation, future source-extension no-go claim,
or program-wide no-go conclusion.

\section{Source Materials}

The AEther-Flow Research Project. (2026, July 2). \emph{Recommendations
implementation plan continue task v15} [Internal implementation plan].

The AEther-Flow Research Project. (2026, July 2). \emph{Narrow source-side
matter-semantics equivalence theorem attempt v1} [Research-control TeX
artifact].

The AEther-Flow Research Project. (2026, July 2). \emph{Source-side
matter-semantics object and certificate manifest v1} [Research-control TeX
artifact].

\end{document}
